\scanpage{22}
\uline{Théorème 2} Le foncteur $K-\Pi_K/\Pi_{\mathbb Q}$ des extensions de t.f. de $Q$ vers les groupoïdes profinis sur $\Pi_{\mathbb Q}$ est \uline{fidèle}.

\margin{En fait, ce théorème n'est pas spécial à $Q$ --- il marcherait sur un corps de base quelconque de car $0$. \unsure{La|car la} démon. montre bien que c'est en fait un énoncé de nature géométrique non arithmétique (\unsure{j'énonce|l'énoncé} pour \unsure{rendre} la démon. facile \unsure{et|une} \unsure{arithmétique|intuitivement} \unsure{pure|courte}).}

Si deux hom. $K\overset{f,g}{\rightrightarrows}K'$ définissent des hom de groupoïdes \inserted{sur $\Pi_{\mathbb Q}$} isomorphes
\diagram{\begin{tikzcd}[ampersand replacement=\&]
\Pi_{K'} \arrow[r,shift left,"f^*"] \arrow[r,shift right,"g^*"'] \arrow[dr,no head,"p'"'] \& \Pi_K \arrow[d,"p"]\\
\& \Pi_{\mathbb Q}
\end{tikzcd}}
(i.e. $\exists$ isomorphisme de foncteurs $\alpha:f^*\simeq g^*$, tel que $\forall\bar\eta'\in\operatorname{Ob}\Pi_{K'}$, l'\deleted{isom}\inserted{carré}
\deleted{commutatif}
\diagram{\begin{tikzcd}[ampersand replacement=\&]
pf^*(\bar\eta') \arrow[r,phantom,"{\overset{p(\alpha)}{\simeq}}" description] \arrow[d,"\wr"'] \& pg^*(\unsure{\eta|\eta'}) \arrow[d,no head,"\wr"]\\
p'(\bar\eta') \arrow[r,phantom,"\simeq" description] \& p'(\bar\eta')
\end{tikzcd}}
soit commutatif) alors $f=g$.

L'hypothèse sur $f,g$ signifie aussi, en termes d'une clôture alg. choisie $\bar K'$ \deleted{\unsure{qui|qu'on} \unsure{peut}}\inserted{donnant via $f,g$ deux clôtures alg. de $K$, $\bar K$, $\widetilde{\bar K}$ de $K$} de $K'$, que \deleted{les deux hom d'exten}\inserted{l'on peut trouver un isomorphisme entre celles-ci (permettant donc d'identifier $E_{\bar K/K}$ et $E_{\widetilde{\bar K}/K}$), de telle façon que les deux hom}
\deleted{\unsure{disons}}
\margin{\unsure{i.e.} induisant l'identité sur les clôtures alg. $\bar Q$ de $Q$ dans $\bar K$ et $\widetilde{\bar K}$, identifiées à celle dans $\bar K'$.}
\[
f^*,g^*:E_{\bar K'/K'}\rightrightarrows E_{\bar K/K}
\]
sont égaux, \deleted{Il \unsure{(Ceci)|C'est} \unsure{s'agit} \unsure{en} \unsure{donc} \unsure{plus}}
\unsure{Ou bien} en termes d'une clôture alg. $\overline{\mathbb Q}$ fixée, de $Q$, en disant que les deux hom $f^*,g^*:\pi_{K'}\rightrightarrows\pi_K$ de groupes profinis extérieurs (avec opérateurs $\Gamma_{\bar Q/Q}$) sont égaux.
\margin{On doit se ramener au cas $k=k'$ --- c'est immédiat en \deleted{\unsure{faisant le chgt de base}}\inserted{\unsure{prenant}} extension finie de $Q$.}
\margin{21}
